\section{Introduction}

Despite LoRA's dominance as the parameter-efficient fine-tuning method of choice, practitioners universally default to rank $r=16$ without principled justification. This one-size-fits-all approach may leave significant performance on the table at scale. As language models grow from billions to hundreds of billions of parameters, the question becomes pressing: should LoRA rank scale with model size?

The surface-level answer might be ``probably,'' but no systematic study has characterized this relationship. Prior work on LoRA has focused on the method itself---low-rank decomposition of weight updates, $\alpha/r$ scaling factors, and which layers to adapt---rather than how optimal rank varies across model scales. This gap leaves practitioners guessing: a 12B model may benefit from rank=64 while a 1B model saturates at rank=16, yet both typically receive the same default configuration.

We address this gap with a systematic empirical study across the Pythia model family (1B to 12B parameters). Our key insight is that \textbf{optimal LoRA rank scales sub-linearly with model size}: $r_{\text{opt}} \propto N^\alpha$ where $\alpha < 1$, with task-dependent values ranging from 0.30 to 0.82. This sub-linear relationship suggests that task-relevant parameter subspaces grow more slowly than overall model capacity---larger models develop more efficient, focused representations requiring proportionally less adaptation.

Beyond the scaling law itself, we uncover two mechanistic findings:

\begin{enumerate}
    \item \textbf{Phase transition in rank sensitivity}: Larger models (12B) exhibit $>$2$\times$ higher sensitivity to rank selection than smaller models (1B), meaning the penalty for suboptimal rank increases with scale.
    \item \textbf{Inverted attention entropy}: Contrary to our initial hypothesis, larger models show \emph{lower} attention entropy (more focused attention patterns), not higher. This suggests larger models develop specialized attention heads rather than broadly distributed patterns.
\end{enumerate}

We also document an important scope limitation: the scaling exponent $\alpha$ is task-dependent. Single-hop QA (SQuAD-v2) yields $\alpha \approx 0.82$ while multi-hop reasoning (HotpotQA) yields $\alpha \approx 0.30$. There is no universal scaling law---task complexity modulates the relationship.

Our contributions are:

\begin{itemize}
    \item \textbf{Empirical scaling law}: First systematic characterization of $r_{\text{opt}}$ vs model size $N$ under controlled conditions
    \item \textbf{Phase transition evidence}: Demonstration that rank sensitivity increases super-linearly with scale
    \item \textbf{Scope calibration}: Honest documentation of task-dependency, preventing overclaiming of a ``universal'' law
    \item \textbf{Mechanistic insight}: Discovery that attention entropy decreases with scale, informing future rank selection methods
\end{itemize}

These findings move LoRA rank selection from ad-hoc defaults toward principled prescriptions, enabling practitioners to configure rank based on model scale and task type.
